\section{课程简介}

这门课是由 MHI，即 Gemeinschaftsvorlesung der Wissenschaftlichen Gesellschaft Montage, Handhabung und Industrierobotik 组织的联合课程。MHI 是由德语区多所大学的教授组成的学术网络，主要研究方向包括装配技术、搬运技术和工业机器人。

后续课程会涉及软件与控制技术、工业机器人、人机交互、传感系统、定位与基于位置的服务、工业数据科学、机器学习与人工智能、仿真与编程技术、工业 4.0 的标准和参考架构，以及增材制造中的数字化等内容。

\subsection{工业4.0的概念}

工业发展的历史通常被划分为\emph{四次工业革命}：
\begin{itemize}
    \item 工业 1.0 指的是以水力和\emph{蒸汽机}为基础的机械化生产。
    \item 工业 2.0 指的是以电力和\emph{流水线}为基础的大规模生产
    \item 工业 3.0 指的是通过电子技术、\emph{信息技术和计算机}实现的自动化生产
    \item 工业 4.0 则指当前阶段的\emph{网络化}、数据化和智能化生产
\end{itemize}

因此，工业 4.0 并不只是“工厂里使用机器人”这么简单。它的核心在于将机器、产品、传感器、软件系统、数据平台和人连接起来，使生产系统能够采集数据、交换信息、分析状态，并在此基础上优化生产过程。工业 4.0 的目标是让生产系统更加灵活、透明、高效，并能够适应更复杂和更个性化的生产需求。

在工业 4.0 中，一个重要概念是\emph{信息物理系统，也就是 Cyber-Physical Systems，简称 CPS}。它指的是物理设备和数字系统高度耦合的系统。例如，一台现代机床不仅仅是机械设备，它还可以通过传感器采集运行数据，通过控制器和网络接口传输数据，并利用软件系统对状态进行监测、分析和优化。这样，现实中的物理过程和数字世界中的信息处理就被连接在一起。

另一个关键概念是 RAMI 4.0，即“工业 4.0 参考架构模型”。它可以理解为描述工业 4.0 系统的三维坐标系。这个模型从三个维度来组织工业 4.0 系统：
\begin{itemize}
    \item 从物理资产到业务层面的不同层级。
    \item 产品或设备从开发、生产、使用到维护的生命周期。
    \item 从单个产品、现场设备、控制设备、工作站、企业到互联世界的层级结构。
\end{itemize}

RAMI 4.0 的作用是帮助人们系统化地理解复杂的工业 4.0 系统。它强调，工业 4.0 不能只从单一技术角度理解，比如只看机器人、传感器或云平台。真正的工业 4.0 涉及物理设备、通信、数据、功能、业务流程、产品生命周期和企业层级之间的整体集成。

工业 4.0 不只是技术问题。除了物联网、人工智能、机器人、云计算和大数据等技术之外，它还涉及开放标准、数字世界与现实世界的集成、网络基础设施、安全、工作组织、教育培训、法律框架和资源效率等问题。

\subsection{物联网 IoT}

IoT 的本质就是让物理世界中的物体具备感知、通信、计算和行动能力，从而与数字世界形成闭环。下面这七点分别对应实现这个闭环所必需的技术前提。
\begin{enumerate}
    \item 泛在通信（Uniquitäre Kommunikation）：物联网设备数量级远超传统互联网终端，而且很多分布在极其分散、甚至难以人工触达的位置。这就要求通信不再局限于Wi-Fi/以太网，而要覆盖低功耗广域网、蜂窝网络、短距无线等多种技术组合，做到随时随地都能连上。
    \item 大数据与云计算（Big Data and Cloud Computing）：海量设备持续产生数据流，单个终端的算力和存储不足以处理。云计算提供弹性的存储与计算资源，大数据技术则负责从这些高维、高频、异构的数据中提取有用信息。
    \item 整体环境感知（Holistische Umwelterfassung）：单一传感器只能提供局部、片面的信息。整体强调的是多模态传感器融合，加上多点位覆盖，才能重建一个足够完整、可用于决策的环境模型。
    \item 所有实体的移动性（Mobilität aller Entitäten）：IoT中的物往往不是静止的。这要求网络架构支持无缝切换、位置感知以及移动场景下的服务连续性。
    \item 能量自主（Hierachiefreie Energieautonomie）：很多IoT节点部署在无法接入电网、也难以频繁更换电池的地方。因此需要不依赖固定基础设施的供能方式，能灵活转换多种能量来源为电能的能量采集技术。这一点直接决定了IoT系统能否实现真正的部署后免维护。
\end{enumerate}
这五点合在一起，构成了从物理感知，到网络传输，再到智能决策的完整技术链条，这也是为什么它们被称为IoT的使能者，而不只是组成部分。少了任一环，系统要么无法规模化部署，要么无法真正自主运行。

